\chapter{Vollständige Distributivgesetze}
\label{b09guide:ch:distribution}

Die folgenden Gesetze verteilen gleichzeitig über eine beliebige
Menge von Auswahlmöglichkeiten. Ihre allgemeine Gültigkeit ist
äquivalent zu AC. Die entsprechenden elementaren Gesetze mit
nur endlich vielen Schritten gehören zur Mengenlehre ohne
Auswahlaxiom und werden hier nicht erneut gesammelt.
Die vollständigen Schnitt- und Vereinigungsgesetze sind in
\cite{collins} als Auswahläquivalenzen behandelt.

\section{Grundmenge und Indexkonventionen}

\begin{definition}[Relativer indizierter Durchschnitt]
Für eine feste Menge $U$ und eine Familie $(E_k)_{k\in K}$
von Teilmengen von $U$ sei
\[
 \bigcap\nolimits^U_{k\in K}E_k
   =\{x\in U\mid\forall k\in K\;(x\in E_k)\}.
\]
Insbesondere bedeutet
$\bigcap^U_{k\in\varnothing}E_k=U$.
Vereinigungen über leere Indexmengen sind weiterhin leer.
\end{definition}

\begin{definition}[Auswahl eines Indexes in jeder Zeile]
Sei $I$ eine Menge, und für jedes $i\in I$ sei $J_i$ eine Menge.
Eine doppelt indizierte Familie $(A_{i,j})$ besitzt genau die
Indizes $(i,j)$ mit $i\in I$ und $j\in J_i$.
Es sei
\[
 D=\prod_{i\in I}J_i.
\]
Ein Element $s\in D$ legt für jede Zeile $i$ genau einen Index
$s(i)\in J_i$ fest. Die Mengen $J_i$ dürfen leer sein;
es wird kein Element von $D$ vorab vorausgesetzt.
\end{definition}

\section{Durchschnitt von Vereinigungen}

\begin{theorem}[Vollständiges Distributivgesetz]
\label{b09guide:thm:distribution}
Seien $U,I,(J_i)_{i\in I}$ beliebig und
$A_{i,j}\subseteq U$ für alle zulässigen $(i,j)$.
Mit $D=\prod_{i\in I}J_i$ gilt unter AC:
\[
 \boxed{\displaystyle
 \bigcap\nolimits^U_{i\in I}\;
       \bigcup_{j\in J_i}A_{i,j}
 =
 \bigcup_{s\in D}\;
       \bigcap\nolimits^U_{i\in I}A_{i,s(i)}.}
\]
Die Aussage für alle diese Mengendaten ist äquivalent zu AC.
\end{theorem}

Das Gesetz beschreibt die gemeinsame Erfüllbarkeit der
Zeilenbedingungen durch eine ganze Auswahl von Indizes.
Für $I=\varnothing$ sind beide Seiten $U$.
Ist ein $J_i$ leer, so sind beide Seiten leer.
Alle Durchschnitte verwenden dieselbe Grundmenge $U$.

\section{Vereinigung von Durchschnitten}

\begin{theorem}[Duales vollständiges Distributivgesetz]
\label{b09guide:thm:dual}
Seien $U,I,(J_i)_{i\in I}$ beliebig und
$A_{i,j}\subseteq U$ für alle zulässigen $(i,j)$.
Mit $D=\prod_{i\in I}J_i$ gilt unter AC:
\[
 \boxed{\displaystyle
 \bigcup_{i\in I}\;
       \bigcap\nolimits^U_{j\in J_i}A_{i,j}
 =
 \bigcap\nolimits^U_{s\in D}\;
       \bigcup_{i\in I}A_{i,s(i)}.}
\]
Auch diese Aussage für alle diese Mengendaten ist äquivalent zu AC.
\end{theorem}

Für $I=\varnothing$ sind beide Seiten leer.
Ist ein $J_i$ leer, so sind beide Seiten $U$.
Die Kennzeichnung des relativen Durchschnitts ist besonders
auf der rechten Seite nötig, wenn $D=\varnothing$.

\section{Produkt von Vereinigungen}

\begin{theorem}[Distributivität des Produkts über Vereinigungen]
\label{b09guide:thm:productdistribution}
Seien $I$, die Mengen $J_i$ und die Mengen $A_{i,j}$ beliebig.
Mit $D=\prod_{i\in I}J_i$ gilt unter AC:
\[
 \boxed{\displaystyle
 \prod_{i\in I}\left(\bigcup_{j\in J_i}A_{i,j}\right)
 =
 \bigcup_{s\in D}\left(\prod_{i\in I}A_{i,s(i)}\right).}
\]
Die Aussage für alle diese Mengendaten ist äquivalent zu AC.
\end{theorem}

Die Gleichheit ist als Gleichheit von Mengen von Funktionsgraphen
gemeint, mit dem Definitionsbereich $I$ für jedes Produktelement.
Eine gesonderte Zielmenge gehört nicht zum Graphen.
Die Auswahl eines Indexes $s(i)$ muß nicht eindeutig sein.
Für $I=\varnothing$ sind beide Seiten $\{\varnothing\}$;
ist ein $J_i$ leer, so sind beide Seiten leer.

\section{Übersicht zum Nachschlagen}

\begin{center}
\small
\begin{tabularx}{\linewidth}{@{}Xl@{}}
\toprule
Gesuchte Aussage & Fundstelle\\
\midrule
Auswahlmengen für disjunkte Familien & Axiom~\ref{b09guide:ax:choice}\\
Auswahlfunktion einer Mengenfamilie & Satz~\ref{b09guide:thm:family}\\
Indizierte Auswahl & Satz~\ref{b09guide:thm:indexed}\\
Auswahl auf allen nichtleeren Teilmengen & Satz~\ref{b09guide:thm:powerset}\\
Nichtleerheit von Produkten & Satz~\ref{b09guide:thm:product}\\
Fortsetzung einer partiellen Auswahl & Satz~\ref{b09guide:thm:extension}\\
Auswahl mit Mengen- und Wertvorgaben & Satz~\ref{b09guide:thm:constraints}\\
Gleichzeitige Auswahlfunktionen & Satz~\ref{b09guide:thm:simultaneous}\\
Disjunkte Verfeinerung einer Überdeckung & Satz~\ref{b09guide:thm:refinement}\\
Durchschnitt von Vereinigungen & Satz~\ref{b09guide:thm:distribution}\\
Vereinigung von Durchschnitten & Satz~\ref{b09guide:thm:dual}\\
Produkt von Vereinigungen & Satz~\ref{b09guide:thm:productdistribution}\\
\bottomrule
\end{tabularx}
\end{center}
